\Needspace{6\baselineskip}
\subsection{Calibration and Test Conditions}
The camera and color sensor are aligned with the original display and pilot, respectively. The temporary mounting uses removable tape and an opaque optical sleeve around the pilot sensor. Camera framing is adjusted through a low-resolution preview; single JPEG captures are then used for recognition. Cropping and a 90-degree rotation are configurable in the image-processing service.

The RGB acquisition uses a 154~ms integration time and 4$\times$ gain. Raw clear-channel values first identify a dark condition. For illuminated samples, the RGB components are divided by the clear reading before comparison with the mounted red and green references. The selected normalized-distance tolerance is 0.07 and the dark cutoff is $C\leq7$. These settings were derived from the recorded calibration observations.

Servo adjustment begins with the mechanism fixed in place: the rest position must release the button, and the press position must register it. Separate parameters are retained for A, B, and C because identical servos can have different mounting orientations. The mounted power-button trial and the implemented command checks are reported separately in the results.

For the flow channel, the firmware converts pulses using its configured factor and measurement interval. A physical calibration would determine $K=N/V_{\mathrm{ref}}$ from an independently measured reference volume. The reported pulse factor is a configuration value, rather than an experimentally fitted coefficient for the heater installation.
